\chapter{谓词逻辑}
\intro{谓词逻辑又称一阶逻辑，是对命题逻辑的扩展。命题逻辑将原子命题视为不可分割的整体，忽略其内部结构；谓词逻辑进一步分析命题的内部成分——个体对象和谓词关系，并引入全称量词 $\forall$ 和存在量词 $\exists$，使推理表达能力大幅增强。}

\section{谓词与个体词}

\subsection{个体词}

\textbf{个体词}是表示所讨论对象的符号，分为两类：
\begin{itemize}
\item \textbf{个体常元}：表示特定对象的符号，用小写字母 $a,b,c,\dots$ 表示。
\item \textbf{个体变元}：表示泛指某类对象的符号，用 $x,y,z,\dots$ 表示。
\end{itemize}

\textbf{论域}（Domain of Discourse）：个体变元的取值范围，记作 $D$。论域可以是有限集或无限集。

\subsection{谓词}

\textbf{谓词}是表示个体性质或个体间关系的符号。谓词本身不是命题，但填入个体词后成为命题。
\begin{itemize}
\item \textbf{一元谓词} $P(x)$：表示个体 $x$ 的性质。
\item \textbf{二元谓词} $R(x,y)$：表示个体 $x$ 与 $y$ 之间的关系。
\item \textbf{$n$ 元谓词} $P(x_1,\dots,x_n)$：表示 $n$ 个个体之间的关系。
\end{itemize}

\begin{center}
\begin{tabular}{c|c|c}
\toprule
谓词符号 & 含义 & 元数 \\
\midrule
$H(x)$ & $x$ 是人 & 一元 \\
$M(x)$ & $x$ 会死 & 一元 \\
$L(x,y)$ & $x$ 爱 $y$ & 二元 \\
$G(x,y,z)$ & $x$ 在 $y$ 和 $z$ 之间 & 三元 \\
\bottomrule
\end{tabular}
\end{center}

命题是 \textbf{0 元谓词}（不含量词和自由变元）。

\subsection{自然语言翻译模板}

\begin{center}
\begin{tabular}{c|c}
\toprule
自然语言 & 谓词公式 \\
\midrule
所有 $P$ 是 $Q$ & $\forall x (P(x) \rightarrow Q(x))$ \\
有些 $P$ 是 $Q$ & $\exists x (P(x) \land Q(x))$ \\
没有 $P$ 是 $Q$ & $\forall x (P(x) \rightarrow \neg Q(x))$ \\
有些 $P$ 不是 $Q$ & $\exists x (P(x) \land \neg Q(x))$ \\
\bottomrule
\end{tabular}
\end{center}

\textbf{关键规律}：全称量词 $\forall$ 后通常跟\textbf{蕴涵} $\rightarrow$，存在量词 $\exists$ 后通常跟\textbf{合取} $\land$。

\begin{example}
将"每个人都爱某个人"翻译为谓词公式。

设论域为全体人类，$L(x,y)$ 表示 "$x$ 爱 $y$"：
\[\forall x \exists y \, L(x,y)\]
注意：$\forall x \exists y \, L(x,y)$（每人都爱某个可能不同的人）与 $\exists y \forall x \, L(x,y)$（存在一个被所有人爱的人）意义不同！
\end{example}

\section{量词}

\subsection{全称量词 $\forall$ 与存在量词 $\exists$}

\begin{mydefinition}{全称量词}
$\forall x \, P(x)$ 读作"对任意 $x$（在论域中），$P(x)$ 成立"。当且仅当论域中的\textbf{每一个}个体都满足 $P$ 时，$\forall x \, P(x)$ 为真。

在有限论域 $D = \{a_1,\dots,a_n\}$ 中：
\[\forall x \, P(x) \equiv P(a_1) \land P(a_2) \land \dots \land P(a_n)\]
\end{mydefinition}

\begin{mydefinition}{存在量词}
$\exists x \, P(x)$ 读作"存在 $x$（在论域中），使得 $P(x)$ 成立"。当且仅当论域中\textbf{至少有一个}个体满足 $P$ 时，$\exists x \, P(x)$ 为真。

在有限论域中：
\[\exists x \, P(x) \equiv P(a_1) \lor P(a_2) \lor \dots \lor P(a_n)\]
\end{mydefinition}

\textbf{唯一存在量词} $\exists!$：$\exists! x \, P(x)$ 表示"存在唯一的 $x$ 使 $P(x)$ 成立"。
\[\exists! x\, P(x) \equiv \exists x\, [P(x) \land \forall y\, (P(y) \rightarrow y = x)]\]

\textbf{量词对偶性}：
\[\neg\forall x P(x) \equiv \exists x \neg P(x),\qquad \neg\exists x P(x) \equiv \forall x \neg P(x)\]

\subsection{辖域、自由变元与约束变元}

\begin{mydefinition}{辖域}
量词的作用范围。在 $\forall x \, P(x) \rightarrow Q(x)$ 中，量词 $\forall x$ 的辖域仅为 $P(x)$，$Q(x)$ 中的 $x$ 是自由变元。在 $\forall x \, (P(x) \rightarrow Q(x))$ 中，辖域是整个 $P(x) \rightarrow Q(x)$。
\end{mydefinition}

\begin{itemize}
\item \textbf{约束变元}：出现在量词辖域中且与该量词使用的变元名相同的变元。
\item \textbf{自由变元}：未受任何量词约束的变元。
\end{itemize}

含自由变元的公式不是命题，其真值依赖于对自由变元的具体赋值。

\begin{example}
区分公式 $\forall x (P(x) \rightarrow \exists y Q(x,y)) \lor R(z)$ 中的自由与约束变元：
\begin{itemize}
\item $\forall x$ 约束了 $P(x)$ 和 $Q(x,y)$ 中的 $x$
\item $\exists y$ 约束了 $Q(x,y)$ 中的 $y$
\item $z$ 是自由变元
\end{itemize}
\end{example}

% --- TikZ: 量词辖域树形图 ---
\begin{figure}[H]
\centering
\begin{tikzpicture}[
    level distance=1.8cm,
    level 1/.style={sibling distance=4cm},
    level 2/.style={sibling distance=2cm},
    dmScopeNode/.style={rectangle,draw=blue!60,fill=blue!5,rounded corners,align=center,font=\small}
]
\node[dmScopeNode] {$\forall x \, (P(x) \rightarrow \exists y \, Q(x,y)) \lor R(z)$}
    child { node[dmScopeNode] {$\forall x$ 辖域\\$P(x) \rightarrow \exists y \, Q(x,y)$}
        child { node[dmScopeNode] {$P(x)$\\$x$ 约束} }
        child { node[dmScopeNode] {$\exists y \, Q(x,y)$\\$x$ 约束, $y$ 约束}
            child { node[dmScopeNode] {$Q(x,y)$\\$x,y$ 约束} }
        }
    }
    child { node[dmScopeNode] {$R(z)$\\$z$ 自由} };
\end{tikzpicture}
\caption{量词辖域树形图——展示约束变元与自由变元的层次关系}
\end{figure}

\subsection{变元的改名规则}

在谓词逻辑中，约束变元的名称可以更改而不改变公式的含义，但必须遵守：
\begin{enumerate}
\item 只对约束变元改名，不能改自由变元。
\item 改名后的变元不能与辖域中已有的其他自由变元同名（避免变量捕获）。
\end{enumerate}

例如：$\forall x \, P(x) \equiv \forall y \, P(y)$（合法）。但 $\forall x \, (P(x) \land Q(y))$ 不能改为 $\forall y \, (P(y) \land Q(y))$（原名自由变元 $y$ 被捕获）。

\section{谓词公式与解释}

\subsection{项与原子公式}

\textbf{项}的递归定义：
\begin{enumerate}
\item 个体常元是项；个体变元是项。
\item 若 $f$ 是 $n$ 元函数符号，$t_1,\dots,t_n$ 是项，则 $f(t_1,\dots,t_n)$ 是项。
\end{enumerate}

\textbf{原子公式}：若 $P$ 是 $n$ 元谓词符号，$t_1,\dots,t_n$ 是项，则 $P(t_1,\dots,t_n)$ 是原子公式。

\subsection{合式公式}

谓词逻辑的\textbf{合式公式}（WFF）递归定义：
\begin{enumerate}
\item 每个原子公式是合式公式。
\item 若 $A$ 是 WFF，则 $(\neg A)$ 是 WFF。
\item 若 $A,B$ 是 WFF，则 $(A \land B)$、$(A \lor B)$、$(A \rightarrow B)$、$(A \leftrightarrow B)$ 是 WFF。
\item 若 $A$ 是 WFF，$x$ 是变元，则 $(\forall x \, A)$ 和 $(\exists x \, A)$ 是 WFF。
\end{enumerate}

\textbf{封闭公式}（Sentence）：不含自由变元的 WFF，具有确定的真值。

\subsection{解释与真值}

谓词公式的\textbf{解释} $\mathcal{I}$ 由以下成分构成：
\begin{enumerate}
\item 论域 $D$：一个非空集合。
\item 常元指派：将每个个体常元 $a$ 映射到 $D$ 中的某个元素。
\item 函数符号指派：将每个 $n$ 元函数符号 $f$ 映射到函数 $f^{\mathcal{I}}: D^n \rightarrow D$。
\item 谓词符号指派：将每个 $n$ 元谓词符号 $P$ 映射到关系 $P^{\mathcal{I}} \subseteq D^n$。
\end{enumerate}

封闭公式在给定解释下按 Tarski 语义递归计算真值。

\begin{example}
设论域 $D=\{1,2,3\}$，$P(x)$:"$x$ 是偶数"（$P^{\mathcal{I}}=\{2\}$），$Q(x)$:"$x>0$"（$Q^{\mathcal{I}}=\{1,2,3\}$）。
\begin{itemize}
\item $\mathcal{I} \models \forall x \, Q(x)$（所有数都 $>0$ ✓）
\item $\mathcal{I} \not\models \forall x \, P(x)$（1 不是偶数 ✗）
\item $\mathcal{I} \models \exists x \, P(x)$（存在偶数 2 ✓）
\end{itemize}
\end{example}

\section{量词的等价与永真式}

\subsection{量词的否定律}

\begin{myformula}
\textbf{量词否定律}
\[
\neg\forall x P(x) \equiv \exists x \neg P(x),\qquad
\neg\exists x P(x) \equiv \forall x \neg P(x)
\]
\end{myformula}

\textbf{直观理解}：$\forall x P(x)$ 为假，意味着并非所有 $x$ 都使 $P(x)$ 成立 → 存在某个 $x$ 使 $P(x)$ 不成立 → $\exists x \neg P(x)$ 为真。

\subsection{量词的分配律}

\begin{center}
\begin{tabular}{c|c|c}
\toprule
编号 & 等价/蕴涵关系 & 说明 \\
\midrule
Q3 & $\forall x (P(x) \land Q(x)) \equiv \forall x P(x) \land \forall x Q(x)$ & $\forall$ 对 $\land$ 可分配 \\
Q4 & $\exists x (P(x) \lor Q(x)) \equiv \exists x P(x) \lor \exists x Q(x)$ & $\exists$ 对 $\lor$ 可分配 \\
Q5 & $\forall x P(x) \lor \forall x Q(x) \models \forall x (P(x) \lor Q(x))$ & 反向不成立 \\
Q6 & $\exists x (P(x) \land Q(x)) \models \exists x P(x) \land \exists x Q(x)$ & 反向不成立 \\
\bottomrule
\end{tabular}
\end{center}

\begin{example}
反例（Q5反向不成立）：论域 $D=\N$，$P(x)$:"$x$ 是偶数"，$Q(x)$:"$x$ 是奇数"。
\begin{itemize}
\item $\forall x (P(x) \lor Q(x))$ 为真（每个自然数非偶即奇）
\item $\forall x P(x) \lor \forall x Q(x)$ 为假（并非所有数都是偶数，也并非所有数都是奇数）
\end{itemize}
\end{example}

\subsection{量词作用域的扩张与收缩}

若公式 $B$ 不含自由变元 $x$，则有以下等价式：

\begin{myformula}
\textbf{量词作用域扩张律}
\[
\begin{aligned}
\forall x (A(x) \land B) &\equiv \forall x A(x) \land B \\
\forall x (A(x) \lor B) &\equiv \forall x A(x) \lor B \\
\exists x (A(x) \land B) &\equiv \exists x A(x) \land B \\
\exists x (A(x) \lor B) &\equiv \exists x A(x) \lor B
\end{aligned}
\]
\end{myformula}

\subsection{量词顺序与互换}

\begin{itemize}
\item $\forall x \forall y \, P(x,y) \equiv \forall y \forall x \, P(x,y)$（同种量词可交换）
\item $\exists x \exists y \, P(x,y) \equiv \exists y \exists x \, P(x,y)$（同种量词可交换）
\item $\exists x \forall y \, P(x,y) \models \forall y \exists x \, P(x,y)$（反向不成立！）
\end{itemize}

\textbf{关键区别}：$\forall y \exists x \, P(x,y)$（对每个 $y$，存在某个可能依赖于 $y$ 的 $x$）\textbf{不能}推出 $\exists x \forall y \, P(x,y)$（存在一个统一的 $x$ 对所有 $y$ 成立）。

\section{前束范式}

\begin{mydefinition}{前束范式}
一个谓词公式称为\textbf{前束范式}，若其具有形式：
\[Q_1 x_1 \, Q_2 x_2 \, \dots \, Q_n x_n \, M\]
其中每个 $Q_i$ 是 $\forall$ 或 $\exists$，$M$ 是不含量词的公式（称为\textbf{母式}），所有量词都在公式的最前面。
\end{mydefinition}

\subsection{求前束范式的四步法}

\begin{enumerate}
\item \textbf{消去 $\rightarrow$ 和 $\leftrightarrow$}：
$A \rightarrow B \equiv \neg A \lor B$；$A \leftrightarrow B \equiv (A \rightarrow B) \land (B \rightarrow A)$

\item \textbf{将否定深入}（德摩根律 + 量词否定律）：
\begin{itemize}
\item $\neg\neg A \equiv A$
\item $\neg(A \land B) \equiv \neg A \lor \neg B$，$\neg(A \lor B) \equiv \neg A \land \neg B$
\item $\neg\forall x A(x) \equiv \exists x \neg A(x)$，$\neg\exists x A(x) \equiv \forall x \neg A(x)$
\end{itemize}

\item \textbf{变元改名}：若不同量词使用了相同的约束变元名，将其中的约束变元改为新名。确保不存在某变元既是约束的又是自由的。

\item \textbf{量词前移}：利用量词作用域扩张律，将量词逐步移到公式前面。
\end{enumerate}

\begin{example}
求 $\forall x P(x) \rightarrow \exists y Q(y)$ 的前束范式。
\[
\begin{aligned}
\forall x P(x) \rightarrow \exists y Q(y)
&\equiv \neg\forall x P(x) \lor \exists y Q(y) &&\text{（消去 $\rightarrow$）} \\
&\equiv \exists x \neg P(x) \lor \exists y Q(y) &&\text{（量词否定律）} \\
&\equiv \exists x \exists y (\neg P(x) \lor Q(y)) &&\text{（量词前移）}
\end{aligned}
\]
\end{example}

% --- TikZ: 前束范式转换流程图 ---
\begin{figure}[H]
\centering
\begin{tikzpicture}[node distance=1.5cm, auto]
\tikzset{
    dmPNF/.style={rectangle,draw=orange!60!black,fill=orange!5!white,rounded corners,minimum width=2.8cm,align=center},
    dmPNFArrow/.style={-{Stealth},thick,orange!60!black}
}
\node[dmPNF] (S1) {原始谓词公式};
\node[dmPNF,below=of S1] (S2) {第一步：消去 $\rightarrow,\leftrightarrow$};
\node[dmPNF,below=of S2] (S3) {第二步：将 $\neg$ 深入};
\node[dmPNF,below=of S3] (S4) {第三步：变元改名};
\node[dmPNF,below=of S4] (S5) {第四步：量词前移};
\node[dmPNF,below=of S5] (S6) {前束范式};
\draw[dmPNFArrow] (S1) -- (S2);
\draw[dmPNFArrow] (S2) -- (S3);
\draw[dmPNFArrow] (S3) -- (S4);
\draw[dmPNFArrow] (S4) -- (S5);
\draw[dmPNFArrow] (S5) -- (S6);
\end{tikzpicture}
\caption{前束范式转换四步法流程}
\end{figure}

\section{谓词逻辑推理}

\subsection{量词的推理规则}

\begin{center}
\begin{tabular}{c|c|c}
\toprule
规则 & 名称 & 形式 \\
\midrule
UI & 全称实例化 & $\forall x P(x) \models P(a)$（$a$ 为任意个体常元）\\
EI & 存在实例化 & $\exists x P(x) \models P(c)$（$c$ 为新常元）\\
UG & 全称泛化 & $P(a)$ 对任意 $a$ 成立 $\models \forall x P(x)$\\
EG & 存在泛化 & $P(a) \models \exists x P(x)$\\
\bottomrule
\end{tabular}
\end{center}

\textbf{关键约束}：
\begin{itemize}
\item EI：$c$ 必须是\textbf{新}常元——此前未在前提、已推出的结论中出现过，且不能出现在最终结论中。
\item UG：$a$ 不能是已被特殊假设（如来自 EI）限制的个体。
\end{itemize}

\subsection{形式证明示例}

\begin{example}
经典三段论：所有人都会死，苏格拉底是人，所以苏格拉底会死。

设 $H(x)$:"$x$ 是人"，$M(x)$:"$x$ 会死"，$s$:"苏格拉底"。证明 $\forall x (H(x) \rightarrow M(x)), H(s) \models M(s)$。
\[
\begin{aligned}
1.\quad &\forall x (H(x) \rightarrow M(x)) &\quad &\text{（前提）} \\
2.\quad &H(s) &\quad &\text{（前提）} \\
3.\quad &H(s) \rightarrow M(s) &\quad &\text{（1, UI）} \\
4.\quad &M(s) &\quad &\text{（2, 3, MP）}
\end{aligned}
\]
\end{example}

\begin{example}
证明 $\forall x (P(x) \rightarrow Q(x)), \exists x P(x) \models \exists x Q(x)$。
\[
\begin{aligned}
1.\quad &\forall x (P(x) \rightarrow Q(x)) &\quad &\text{（前提）} \\
2.\quad &\exists x P(x) &\quad &\text{（前提）} \\
3.\quad &P(c) &\quad &\text{（2, EI：引入新常元 $c$）} \\
4.\quad &P(c) \rightarrow Q(c) &\quad &\text{（1, UI）} \\
5.\quad &Q(c) &\quad &\text{（3, 4, MP）} \\
6.\quad &\exists x Q(x) &\quad &\text{（5, EG）}
\end{aligned}
\]
\end{example}

\begin{example}
证明 $\exists x \forall y \, L(x,y) \models \forall y \exists x \, L(x,y)$。

注意：其逆（$\forall y \exists x L(x,y) \models \exists x \forall y L(x,y)$）不成立。直观：每人爱某个人（依赖 $y$ 的 $x$）$\not\Rightarrow$ 存在一个被所有人爱的人（统一的 $x$）。
\end{example}

% --- TikZ: 推理规则应用示意图 ---
\begin{figure}[H]
\centering
\begin{tikzpicture}[node distance=1.5cm, auto, font=\small]
\tikzset{
    dmRule/.style={rectangle,draw=blue!60,fill=blue!10,rounded corners,minimum width=2.4cm,align=center},
    dmRuleArrow/.style={-{Stealth},thick,blue!40}
}
\node[dmRule] (UI) {UI\\全称实例化};
\node[dmRule,right=of UI,xshift=1cm] (EI) {EI\\存在实例化};
\node[dmRule,right=of EI,xshift=1cm] (UG) {UG\\全称泛化};
\node[dmRule,right=of UG,xshift=1cm] (EG) {EG\\存在泛化};
\node[dmRule,above=of UI,yshift=-0.5cm] (Forall) {$\forall x P(x)$};
\node[dmRule,above=of EI,yshift=-0.5cm] (Exists) {$\exists x P(x)$};
\node[dmRule,below=of UI,yshift=0.5cm] (Pa) {$P(a)$ 任意 $a$};
\node[dmRule,below=of EI,yshift=0.5cm] (Pc) {$P(c)$ 新 $c$};
\node[dmRule,below=of UG,yshift=0.5cm] (Forall2) {$\forall x P(x)$};
\node[dmRule,below=of EG,yshift=0.5cm] (Exists2) {$\exists x P(x)$};
\draw[dmRuleArrow] (Forall) -- (UI);
\draw[dmRuleArrow] (UI) -- (Pa);
\draw[dmRuleArrow] (Exists) -- (EI);
\draw[dmRuleArrow] (EI) -- (Pc);
\draw[dmRuleArrow] (Pa) -- (UG);
\draw[dmRuleArrow] (UG) -- (Forall2);
\draw[dmRuleArrow] (Pc) -- (EG);
\draw[dmRuleArrow] (EG) -- (Exists2);
\end{tikzpicture}
\caption{谓词逻辑四大约束规则关系图——UI与UG配对，EI与EG配对}
\end{figure}

\subsection{常见推理错误}

\begin{center}
\begin{tabular}{c|c}
\toprule
错误 & 描述 \\
\midrule
逆命题错误 & 从 $P(x) \rightarrow Q(x)$ 推出 $Q(x) \rightarrow P(x)$ \\
量词跳跃 & 从 $\forall y \exists x P(x,y)$ 错误推出 $\exists x \forall y P(x,y)$ \\
EI使用不当 & 引入的常元在前提中已出现 \\
UG使用不当 & 对由EI引入的特定个体使用全称泛化 \\
\bottomrule
\end{tabular}
\end{center}

\begin{mycode}{有限论域量词评估}
\begin{lstlisting}[language=Python]
from typing import Callable, TypeVar, List

T = TypeVar('T')

def forall(domain: List[T], pred: Callable[[T], bool]) -> bool:
    """评估 ∀x P(x) 在有限论域上的真值"""
    return all(pred(x) for x in domain)

def exists(domain: List[T], pred: Callable[[T], bool]) -> bool:
    """评估 ∃x P(x) 在有限论域上的真值"""
    return any(pred(x) for x in domain)

def exists_unique(domain: List[T], pred: Callable[[T], bool]) -> bool:
    """评估 ∃!x P(x)"""
    return sum(1 for x in domain if pred(x)) == 1

# 示例：论域 D = {1,2,3,4,5}
D = [1, 2, 3, 4, 5]
is_even = lambda x: x % 2 == 0
is_positive = lambda x: x > 0

print(f"∀x (x>0) = {forall(D, is_positive)}")  # True
print(f"∀x (x是偶数) = {forall(D, is_even)}")  # False
print(f"∃x (x是偶数) = {exists(D, is_even)}")  # True

# 验证量词否定律：¬∀x Q(x) ≡ ∃x ¬Q(x)
lhs = not forall(D, is_even)
rhs = exists(D, lambda x: not is_even(x))
print(f"¬∀x Q(x) ≡ ∃x ¬Q(x): {lhs == rhs}")  # True

# 验证 Q6 反向不成立：(∃xP ∧ ∃xQ) → ∃x(P∧Q)
P = lambda x: x % 2 == 0
Q = lambda x: x % 3 == 0
fwd = exists(D, lambda x: P(x) and Q(x))
bwd = exists(D, P) and exists(D, Q)
print(f"反向不成立验证: {not bwd or fwd}")  # False
\end{lstlisting}
\end{mycode}
